\documentclass[11pt]{article}

% --------- BASIC PACKAGES ----------
\usepackage[left=0.68in,top=0.6in,right=0.68in,bottom=0.6in]{geometry}
\usepackage{enumitem}
\usepackage[hidelinks]{hyperref}
\usepackage{titlesec}
\usepackage{array}
\usepackage{setspace}

\usepackage[T1]{fontenc}
\usepackage{newtxtext,newtxmath} % Times-like
\input{glyphtounicode}
\pdfgentounicode=1

\setstretch{1.0}
\pagenumbering{gobble}

% ---------- SECTION FORMATTING ----------
\titleformat{\section}{\Large\scshape}{}{4em}{}[\vspace{-0.7em}\rule{\linewidth}{0.4pt}\vspace{-0.4em}]
\titlespacing*{\section}{0pt}{1em}{0.4em}

\newenvironment{cvrole}[5]{%
  \par\noindent\textbf{\large #1} | \textit{\small #2, #3}\hfill \textit{\small #4 -- #5}\par
  \begin{itemize}[leftmargin=2em, itemsep=-0.3em, topsep=-0.5em]%
}{%
  \end{itemize}\par\noindent\mbox{}\par\vspace{0.1em}%
}

% ---------- HEADER ----------
\begin{document}
\pagestyle{empty}

\begin{center}
    {\LARGE \scshape \textbf{Nishal Stanislaus Silva}}{\large \textit {, PhD}}
    \\[1pt]
    {\small Machine Learning Research Engineer | Audio and Acoustic Machine Learning, Edge Inference, Benchmark Design}
    \\[1pt]
    {\footnotesize\textit{Canadian Permanent Resident, Eligible to work in Canada without sponsorship}}
    \\[1pt]
    {\small
    \href{mailto:nishal.silva@hotmail.com}{nishal.silva@hotmail.com}
    \,\,|\,\, +1 613 200 2030
    \,\,|\,\, Toronto, ON M6P 1Z2
    \,\,|\,\, \href{https://nishal.xyz}{nishal.xyz}
    \,\,|\,\, \href{https://www.linkedin.com/in/nishal-silva}{linkedin.com/in/nishal-silva}}
\end{center}

% =====================================================
\section*{Professional Summary}

Machine learning research engineer with a PhD and over ten years across industry and academia, focused on deep learning for audio under hard latency and compute constraints. My doctoral and postdoctoral work delivered real-time acoustic pattern detection on edge-class hardware: F1 of 0.76 at 14\,ms inference on a Raspberry Pi 4, a $74\times$ speedup over DSP baselines. I am well experienced in pre-model evaluation, comprehensive benchmarks, ablations, and I have extensive experience in conducting comprehensive and holistic model evaluations. I have released four public audio benchmark datasets with explicit label ontologies. I have published my work at the in Journal of Audio Engineering Society, multiple IEEE venues, AudioMostly, and Digital Audio Effects conference; My scientific service includes review for IEEE Access, Journal of the National Science Foundation of Sri Lanka, program committee for the IEEE Internet of Sounds Symposium.

% =====================================================
\section*{Core Competencies}

\textbf{Audio and Acoustic Machine Learning:} Sound Event Detection, Audio Classification, Real-Time Audio Pattern Detection, Audio Representation Learning, Feature Front-End Design, Spectrograms and Time--Frequency Trade-offs, Mel and MFCC Representations, Onset Detection, Multimodal Audio and Sensor Learning\\
\textbf{Edge and Constrained Inference:} Embedded and Edge Deployment, Real-Time Streaming Inference, Latency, Power and Compute Profiling, Designing to a Fixed Compute Budget, C++ Real-Time Inference Engines, ONNX\\
\textbf{Evaluation and Research Method:} Benchmark Design, Frozen Evaluation Protocols, Ablation Studies, Cross-Site and Cross-Device Generalisation, Dataset Curation and Label Ontology Design, Weak and Scarce Annotation Regimes, Classical Statistics and Model Validation, Peer-Reviewed Publication

% =====================================================
\section*{Technical Stack}

\textbf{Languages:} Python, C++, C, MATLAB, Linux Shell Scripting, JavaScript\\
\textbf{Machine Learning and Deep Learning:} PyTorch, TensorFlow, Keras, scikit-learn, HuggingFace Transformers, ONNX, NumPy, SciPy, Pandas\\
\textbf{Audio and Signal Processing:} Librosa, JUCE, Elk Audio OS, MIDI, Real-Time DSP, Custom C++ Feature Extraction\\
\textbf{Edge and Systems:} Raspberry Pi, Embedded Linux, Real-Time Audio I/O, IoT Sensor Fleets, Distributed Field Deployment\\
\textbf{Development and MLOps:} Docker, Git, Linux/UNIX, Pytest, CI/CD, Jenkins, Apache Airflow, AWS

% =====================================================
\section*{Education}
\textbf{Ph.D - Information Engineering and Computer Science} \textit{\small{ - University of Trento, Italy}} \hfill \textit{Jan 2025}\\[2pt]
\noindent\textbf{M.Sc - Telecommunication and Electronic Engineering} \textit{\small{ - Sheffield Hallam University, UK}} \hfill \textit{Aug 2018}\\[2pt]
\noindent\textbf{B.Eng (Hons) - Electronic Engineering} \textit{\small{ - Sheffield Hallam University, UK}} \hfill \textit{Mar 2014}

\section*{Portfolio and Demos \footnotesize {\normalfont{(full portfolio available at: \href{https://nishal.xyz/research}{\textit{nishal.xyz/research}}.})}}
\begin{itemize}[leftmargin=1.2em, itemsep=-0.25em]
    \item Nebula, a comprehensive C++ audio feature extraction library: \href{https://github.com/ns2max/nebula}{\textit{https://github.com/ns2max/nebula}}
    \item Hot Licks, a real-time audio pattern detection system running on embedded hardware: \href{https://youtu.be/gkOqfOT8zXM}{\textit{https://youtu.be/gkOqfOT8zXM}}
    \item Multisensory concert enabled by real-time on-device audio inference: \href{https://youtu.be/axEHkdnxFB8}{\textit{https://youtu.be/axEHkdnxFB8}}
    \item A connected, synchronized Internet of Musical Things Performance Ecosystem: \href{https://youtu.be/jfm5L3DYIcc}{\textit{https://youtu.be/jfm5L3DYIcc}}

\end{itemize}


\newpage
% =====================================================
\section*{Professional Experience}

\begin{cvrole}{Independent Machine Learning Researcher}{Self-Directed Research}{Toronto, Canada}{Jan 2026}{Present}
    \item Currently researching foundation-model-based music generation and synthetic variation generation for training under low labelled data.
    \item Researching on deep learning-based methods to qualitatively evaluate pronunciation for speech applications.
    \item Released Nebula, an open-source C++ audio feature extraction library designed for constrained hardware.
    \item 2 manuscripts on edge-to-server offloading and real-time pattern detection accepted for publication.
    % \item Curating and releasing public audio benchmark datasets with explicit label ontologies and documented collection protocols, extending a four-dataset suite used for real-time detection research.
    % \item Preparing a controlled comparison of audio against symbolic input representations for embedded real-time drum pattern recognition, currently in review at the Journal of the Audio Engineering Society.
\end{cvrole}

\begin{cvrole}{Postdoctoral Researcher}{University of Trento}{Trento, Italy}{Jan 2025}{Dec 2025}
    \item Owned full research pipelines within the EU-funded MUSMET project on streaming audio and sensor data: collection, feature engineering, training and evaluation, and deployment to edge hardware with runtime monitoring.
    \item Delivered real-time audio pattern detection at 14\,ms inference on a Raspberry Pi 4 at F1 0.76, experimentally demonstrated that recurrent models over raw polyphonic audio stay viable inside a hard real-time budget.
    \item Designed frozen benchmark protocols with baselines and ablations to quantify architecture and feature trade-offs under latency and compute constraints.
    \item Produced three peer-reviewed publications in 2025 (IEEE Internet of Sounds, IEEE I3DA) and presented demonstrations at international research and industry events.
\end{cvrole}

\begin{cvrole}{Visiting Researcher}{McGill University}{Montreal, Canada}{May 2024}{Aug 2024}
    \item Prototyped diffusion-based and VAE-based synthetic audio generation methods to improve robustness where labelled data was scarce.
    \item Developed a real-time multimodal detection system over combined audio and inertial streams; profiled compute and power draw to assess feasibility for inference on battery-powered devices.
\end{cvrole}

\begin{cvrole}{Visiting Researcher}{University of Visual and Performing Arts}{Colombo 02, Sri Lanka}{Jan 2024}{Mar 2024}
    \item Ran structured evaluations of machine learning driven interactive systems with domain practitioners, combining quantitative metrics with expert judgement to produce usability and adoption recommendations.
    \item Worked directly with domain experts to define what counted as correct behaviour; translated practitioner judgements into evaluation criteria to measure models.
\end{cvrole}

\begin{cvrole}{Technical Consultant}{Forestpin (Pvt) Ltd}{Colombo 05, Sri Lanka}{Aug 2020}{Feb 2021}
    \item Designed and deployed statistical and machine learning pipelines for anomaly detection over large structured enterprise datasets, supporting financial forensics and compliance monitoring.
    \item Built Python backend scoring services to flag anomalous records for expert triage.
    % , tuned so limited analyst review time went to genuine exceptions rather than false positives.
    \item Translated regulatory and compliance requirements from non-technical stakeholders into detection specifications with measurable acceptance criteria.
\end{cvrole}

\begin{cvrole}{Research Engineer}{MAS Holdings (Pvt) Ltd}{Colombo, Sri Lanka}{Jan 2015}{Jul 2020}
    \item Owned end-to-end development of machine learning and computer vision systems for industrial manufacturing. 
    \item Integrated camera, sensor, and IoT streams into live production workflows; deployed systems that improved operational efficiency by up to 300\%.
    \item Developed and deployed C++ and Python real-time inference engines on embedded and IoT hardware, optimized for throughput, latency, and reliability within fixed hardware budgets.
    \item Maintained deployed models across distributed sites with differing sensors and operating conditions; Addressed for the cross-site and cross-device drift that appeared in the field but not in development.
    \item Deployed human-in-the-loop inspection models balancing throughput against operational risk; mentored engineers and established code review and CI/CD practice to scale adoption across teams.
\end{cvrole}

% =====================================================
\section*{Selected Publications \footnotesize {\normalfont{(full list of publications available at: \href{https://nishal.xyz/publications}{\textit{nishal.xyz/publications}}.})}}
\begin{itemize}[leftmargin=1.2em, itemsep=-0.25em]
    \item Silva, N. and Turchet, L. \textit{Music Information Retrieval as a Service (MIRaaS): Latency Characterization of an Edge-to-Server Architecture for Near-Real-Time Audio Analysis}. IEEE International Symposium on the Internet of Sounds, 2026. \textit{(accepted)}
    \item Silva, N. and Turchet, L. \textit{Audio vs. MIDI in Embedded Real-Time Drum Pattern Recognition}. Journal of the Audio Engineering Society, 2026. \textit{(accepted)}
    \item Silva, N. and Turchet, L. \textit{Real-Time Audio Pattern Detection for Smart Musical Instruments}. Journal of the Audio Engineering Society, Mar. 2026. 
    \item Silva, N., Wanderley, M. and Turchet, L. \textit{Melody and Motion: Integrating Guitar Gesture Detection with Musical Patterns for Extended Control in Live Performance and Metaverse Applications}. IEEE International Symposium on the Internet of Sounds, Oct. 2025.
    \item Silva, N., Boem, A. and Turchet, L. \textit{Interactive IoMusT-Based Concerts: Real-Time Pattern Recognition and Audience Experience}. IEEE International Conference on Immersive and 3D Audio, Sep. 2025.
    \item Silva, N. and Turchet, L. \textit{User-Centered Evaluation of Smart Musical Instruments with Embedded Real-Time Pattern Detection}. IEEE International Conference on Immersive and 3D Audio, Sep. 2025.
    \item Silva, N. and Turchet, L. \textit{Real-Time Pattern Recognition of Symbolic Monophonic Music}. 19th International Audio Mostly Conference, Sep. 2024.
    \item Silva, N. and Turchet, L. \textit{A Structural Similarity Index Based Method to Detect Symbolic Monophonic Patterns in Real-Time}. 25th International Conference on Digital Audio Effects (DAFx), Sep. 2022.
    \item Silva, N., Fiscione, C. and Turchet, L. \textit{Towards Real-Time Detection of Symbolic Musical Patterns: Probabilistic vs. Deterministic Methods}. 27th FRUCT Conference, Apr. 2021.
    \item Silva, N., Fiscione, C. and Weeraddana, P.C. \textit{On Musical Onset Detection via the S-Transform}. 52nd Asilomar Conference on Signals, Systems, and Computers, Oct. 2018.
\end{itemize}

\section*{Public Benchmark Datasets \footnotesize {\normalfont{(7,000+ recordings, 70 contributors, released on Zenodo)}}}
\begin{itemize}[leftmargin=1.2em, itemsep=-0.25em]
    \item \textbf{DoMP}, 4,000 recordings, symbolic monophonic patterns with expressive variations (\textit{DOI: \href{https://doi.org/10.5281/zenodo.10818617}{\textit{10818617}}})
    \item \textbf{DoPP}, 2,000 recordings, polyphonic audio patterns with expressive variations (\textit{DOI: \href{https://doi.org/10.5281/zenodo.14497998}{\textit{14497998}}})
    \item \textbf{DoDP}, 2,000 recordings, percussive audio, rhythmic pattern detection (\textit{DOI: \href{https://doi.org/10.5281/zenodo.14497974}{\textit{14497974}}})
    \item \textbf{DoDP2}, 2,000 recordings, percussive audio, extended ontology and multi-performer coverage (\textit{DOI: \href{https://doi.org/10.5281/zenodo.18395007}{\textit{18395007}}})
\end{itemize}

\section*{Scientific Committee and Peer Review}
\begin{itemize}[leftmargin=1.2em, itemsep=-0.25em]
    \item \textbf{Reviewer:} IEEE Access, Journal of the National Science Foundation Sri Lanka, Colloquio di Informatica Musicale.
    \item \textbf{Programme Committee:} IEEE Internet of Sounds Symposium, IEEE Conference on Immersive and 3D Audio.
\end{itemize}

\section*{Selected Projects}
\begin{itemize}[leftmargin=1.2em, itemsep=-0.25em]
    \item \textbf{Real-Time Audio Detection on Edge Hardware:} Recurrent detection model over raw polyphonic audio reaching F1 0.76 at 14\,ms inference on a Raspberry Pi 4, $74\times$ faster than the baseline it replaced.
    \item \textbf{Nebula, C++ Audio Feature Extraction Library:} Open-source library implementing spectral, temporal, and cepstral feature front-ends, written to run inside the audio callback on constrained hardware without allocation.
    \item \textbf{Training-Free Detection via Structural Similarity:} Repurposed the structural similarity index from computer vision to match time--frequency structure directly, achieving 95\% detection with zero training data.
    \item \textbf{Multimodal Audio and Inertial Detection:} Combined audio and IMU streams for real-time gesture recognition with McGill University; developed a prototype instrument. IEEE Internet of Sounds, 2025.
\end{itemize}

\vspace{-1em}
\section*{Highlights and Recognition}
\begin{itemize}[leftmargin=1.2em, itemsep=-0.25em]
    \item \textbf{MIDI Innovation Awards:} Finalist, Sep 2023, for a real-time embedded audio pattern detection system.
    \item \textbf{GMOA Recognition (Sri Lanka):} Computer vision system for COVID-19 patient vitals monitoring, deployed across hospital wards during the pandemic.
\end{itemize}

\end{document}